\section{Dataset, outcomes, and preprocessing}
The project uses the UCI Myocardial Infarction Complications dataset, a retrospective clinical dataset containing admission characteristics and outcomes \cite{ucimi}. The source file and variable dictionary are retained in the project; Stage~1 outputs supply the development/test partitions, predictor tables, and role documentation. The available project summary records 1,700 admissions: 1,360 development records and 340 held-out records. Development mortality prevalence is 15.96\%; test mortality prevalence is 15.88\% as recorded in the split summary and aggregate T07 report.

The targets are treated as separate binary tasks. Death is defined as $\mathbb{1}(\texttt{LET\_IS}>0)$ and the held-out evaluator checks consistency with that definition. Predictors exclude IDs, all outcome columns, and the documented post-admission fields. The feature matrices are the existing Stage~1 tables; downstream stages do not reconstruct the split.

Preprocessing is estimator-specific. Random Forest pipelines perform fold-local median numeric and mode categorical imputation followed by one-hot encoding for nominal features. The selected CatBoost model uses native categorical features and its supported numeric missing-value behavior, with categorical missingness represented by the training pipeline's sentinel convention. No class weighting or resampling was used for the frozen finalists; probabilities therefore reflect natural development prevalence. Learned preprocessing is fitted within training folds, not before cross-validation. Full detail and feature roles are in \texttt{mi\_stage1\_eda.py}, \texttt{mi\_column\_dictionary.csv}, and the Stage~1/5 audit outputs.

\subsection{Split and validation design}
Stage~1 fixes a stratified 80/20 split using random seed 42, stratified on \texttt{DEATH\_BINARY}. This yielded 1,360 development and 340 test records. Development model comparisons use five stratified folds with seed 42 and target-specific eligible rows. Average precision (AP, equivalent to non-interpolated area under the precision-recall curve in these reports) is the primary ranking/tuning criterion; AUROC, Brier score, and log loss are secondary. The Stage~6 report checks fold signatures against preceding stages. Test labels were not used for model or threshold selection. The final frozen model, calibrator, and threshold were fixed in T06 before the single T07 test evaluation.
